\section{Introduction}\label{sec:intro}

\subsection{Silent transitions}

A language model can produce more pages of mathematical text in an afternoon than a careful reader can check in a week. The errors
that matter in such a body of text are usually not the obvious ones. A sign error or a wrongly expanded bracket tends to surface at the
next calculation that uses it. Much harder to see are steps that leave no trace in the prose at all. We call them \emph{silent
transitions}, and three kinds recur:
\begin{itemize}[leftmargin=1.5em]
\item \emph{Promotion.} A statement that holds only under a stated condition, or that was only checked numerically, is cited a few
pages later as if it were a theorem, and further results are built on it.
\item \emph{An input read as a result.} A value that was taken from the literature, fitted to data or simply chosen reappears later
as a prediction of the theory.
\item \emph{Frame mixing.} Two quantities are compared or combined although they belong to different frames of description (different
observers, conventions, units, or model assumptions), and no map between those frames has been stated.
\end{itemize}
None of these is an arithmetic error, and each is a natural thing for a language model to do. A model continues a text so that it reads
coherently. A condition that was mentioned twenty paragraphs earlier contributes little to the coherence of the next sentence, so the
continuation smooths it away. This is not a rare malfunction but the expected behaviour of a system trained to continue text fluently.
Human authors do the same under time pressure; language models simply do it at scale.

We call the underlying problem \emph{premise drift}: over a long record, the premises on which a statement actually rests
drift away from the premises the text claims for it. Rewording does not reveal drift, because rewording is exactly how it happens.
What is needed is a structure outside the prose that records what each statement rests on and that cannot be changed by
rephrasing.

\subsection{What this article offers}

This article describes such a structure, called QAM, and the software that checks it. QAM grew out of a long, language-model-assisted
research project in mathematical physics~\cite{repo}, but nothing in it depends on that project, and this article assumes no
knowledge of it. Everything needed to understand the system, including the proofs of its properties, is contained here.

The contributions are:
\begin{enumerate}[leftmargin=1.5em]
\item A \emph{model} of research records (Section~\ref{sec:model}): an append-only sequence of typed entries, each with a frame, a
status, a set of inputs and a list of premises, and four rules that every entry must satisfy.
\item \emph{Proofs} of what the rules guarantee for whole records (Section~\ref{sec:properties}): no conclusion is ever stronger than
anything it rests on, however long the chain; inputs are never lost and never invented; frames change only through declared maps;
and a hash chain binds every conclusion to the exact versions of its premises. The core statements are also machine-checked in Lean~4
(Appendix~\ref{app:lean}).
\item \emph{Frames as an interface} (Section~\ref{sec:frames}): ordinary formulas are written in ordinary notation, but each one is
registered under a declared frame and interpretation, and moving it to another frame requires an explicit, declared step.
\item \emph{Tools} that check a record independently of one another: a \LaTeX{} package that refuses to typeset a record that breaks
the rules, a Lean program that checks the exported record against the same rules, and a script that recomputes all hashes (Section~\ref{sec:latex}).
\item An \emph{execution layer} (Sections~\ref{sec:requests}--\ref{sec:tape}): calculations are submitted as requests to an executor
that answers with one of three outcomes and stores a hashed receipt; a premise graph answers what a result rests on and what
depends on a premise; and a small tape machine is executed both in software and in two hardware descriptions as a cross-check.
\end{enumerate}

\subsection{Two design principles}

Two old ideas from computing shaped the design, and we state them because they explain decisions that would otherwise look
arbitrary.

The first is Konrad Zuse's separation of the arithmetic unit from the plan it executes~\cite{rojas97,bauer72}. Zuse's Z3 carried out
whatever the punched film prescribed; his Plankalk\"ul described computations independently of any machine, with separate classes of
variables for inputs and results. QAM adopts the same division of roles. The language model, like a computer algebra system or a
simulator, is the \emph{arithmetic unit}: it may formulate, transform and calculate. Which statements exist, what status they have and
what they rest on is the \emph{plan}, and the plan is a record that separate checkers read. Inputs and derived results are different
kinds of entry, just as Zuse's input and result variables were.

The second is Alan Turing's insight that not every question can be decided mechanically~\cite{turing36}. A checking system that knows
only ``valid'' and ``invalid'' must therefore either reject everything it cannot decide or quietly accept it. QAM gives the undecided
case its own outcome, the \emph{guard}: an open condition that is stored, hashed and inherited by everything built on it, rather than
smoothed away.

\subsection{How to read this article}

Section~\ref{sec:example} introduces a small worked example from electrodynamics that is used throughout. Sections~\ref{sec:model}
and~\ref{sec:properties} give the model and prove its properties. Sections~\ref{sec:frames} to~\ref{sec:llm} describe the tools, each
introduced on the same example. Section~\ref{sec:related} relates QAM to proof assistants, truth maintenance and provenance, and
collects the limitations. Commands for reproducing every result are collected in Appendix~\ref{app:repro}; the text itself avoids them
where it can.
